\section{Simulation studies}\label{sec:simulations}

The simulations examine conditional laws of quantile-valued outcomes under selection on observed covariates. We compare WCF with DRF and Causal-DRF in location and shape-effect designs, symmetric nonlinear designs, and a structural-zero mixture. Each replication uses the same training sample and independent test covariates across methods. The principal sample sizes are $n=500$ and $n=1000$, with ten replications per setting. Each test sample contains 1000 units. The moderator $V$ divides $X_1$ into four intervals with boundaries $-0.5$, $0$, and $0.5$.

\subsection{Location and shape-effect designs}\label{sec:income-design}

Covariates are independent $\mathrm{Unif}[-1,1]$ variables. For these designs $d=6$, and $A\mid X\sim\mathrm{Bernoulli}\{e(X)\}$. Potential quantile vectors follow
\begin{equation}\label{eq:dgp}
 Q_k^a=m_a(X)+\xi_a+\exp\{s_a(X)+\eta_a\}\psi\{z_k;\gamma_a(X)\},
 \qquad z_k=\Phi^{-1}\{(k-1/2)/K\},
\end{equation}
where
\[
 \psi(z;\gamma)=z+\gamma\left\{\frac{z^2-1}{2}+\frac{z^3-3z}{6}\right\}.
\]
Within each arm, $\xi_a\sim N(0,0.20^2)$ and $\eta_a\sim N(0,0.12^2)$ are independent of one another and of $X$. Assignment depends only on $X$. The function $\psi$ is increasing for $0\leq\gamma<1$, since its derivative is $1-\gamma+\gamma(z+1)^2/2$.

Write
\begin{align*}
 b(x)&=0.55x_2+0.35x_3-0.45x_4+0.30x_5,\\
 s(x)&=0.16+0.08x_6-0.04x_2,\qquad
 g(x)=0.55+0.10x_4-0.08x_2,
\end{align*}
and let $C_{l,u}(t)=\min\{u,\max(l,t)\}$ and $\sigma(t)=(1+e^{-t})^{-1}$. The four designs are
\begin{align*}
\mathrm{IC0}:\quad &m_a=b,\quad s_a=s,\quad \gamma_a=C_{.05,.85}(g),\\
\mathrm{IC1}:\quad &m_a=b+a(0.20+0.15x_4),\quad s_a=s,\quad
 \gamma_a=C_{.05,.85}\{g-a(0.16+0.06x_4)\},\\
\mathrm{IC2}:\quad &m_a=b+a(0.10+0.05x_3),\quad s_a=s-0.12a,\quad
 \gamma_a=C_{.05,.85}(g-0.08a),\\
\mathrm{IC3}:\quad &(m_a,s_a,\gamma_a)\text{ as in IC1}.
\end{align*}
IC0 and IC1 use $e(x)=C_{.10,.90}\{\sigma(1.4x_4-1.1x_5+0.4x_6)\}$. IC2 uses $C_{.05,.95}\{\sigma(1.8x_4-0.9x_5)\}$, and IC3 uses $C_{.01,.99}[\sigma\{3.5(x_4-0.9x_5)\}]$. Thus IC0 has no treatment effect, IC1 changes location and shape, IC2 additionally changes scale, and IC3 reduces overlap while keeping IC1's outcome surfaces. The reference vector is
\[
 q_{\star,k}=1.10+0.60\left[z_k+0.30\left\{\frac{z_k^2-1}{2}+\frac{z_k^3-3z_k}{6}\right\}\right].
\]
These equations describe a general location, scale, and shape construction. A regional income distribution is one possible interpretation, with treatment changing its location or dispersion and adoption depending on regional characteristics. The simulated curves are not constrained to be nonnegative, so this interpretation is illustrative rather than a calibrated model of income or expenditure.

\subsection{Symmetric and nonlinear designs}\label{sec:nonlinear-design}

The additional designs use $d=5$ independent uniform covariates and equation~\eqref{eq:dgp}, with reference $q_{\star,k}=z_k$. Define
\[
 f(x)=0.6\sin(\pi x_1)+0.4x_2x_3,\qquad
 e_m(x)=C_{.10,.90}[\sigma\{0.8(x_1+0.5x_2)\}].
\]
The conditional quantile contours are symmetric when $\gamma_a=0$, even though their locations and scales vary across units. D0 uses $m_a=f+0.8ax_1$, $s_a=0.20x_4+0.25ax_2$, $\gamma_a=0$, and zero latent shocks. D2 is a null-effect design with $m_a=f$, $s_a=0.20x_4$, $\gamma_a=0$, and independent normal shocks with standard deviations $(0.35,0.20)$.

D5 has equal mean quantile curves but different conditional laws. Its location is $m_a=f$ and $\gamma_a=0$. The control arm has $s_0=0.20x_4$ and shock standard deviations $(0.40,0.45)$; the treated arm has $s_1=0.20x_4+0.45^2/2$ and shock standard deviations $(0.15,0)$. The log-scale adjustment equates the expected scales. D6 uses $m_a=f+0.7ax_1$, $s_a=0.20x_4$, and $\gamma_a=0$, with $\eta_a\sim N(0,0.15^2)$ and
\[
 \xi_a\sim\tfrac12 N(-1.5,0.25^2)+\tfrac12 N(1.5,0.25^2).
\]
This creates a multimodal law of quantile curves while every realized curve retains a symmetric inner distribution. D7 changes shape: $m_a=f$, $s_a=0.20x_4$, $\gamma_a=0.20+0.10x_3+a(0.30+0.10x_1)$, and shock standard deviations $(0.35,0.20)$. These five designs use $e_m$.

D8 combines a smaller treatment effect with stronger confounding. It uses
\begin{align*}
 m_a(x)&=1.2\sin(\pi x_1)+0.9x_2x_3+a(0.12x_1+0.06),\\
 s_a(x)&=0.20x_4+0.05ax_2,\qquad \gamma_a(x)=0,\\
 e(x)&=C_{.05,.95}[\sigma\{2.5(x_1+0.7x_2-0.5x_3)\}],
\end{align*}
with shock standard deviations $(0.35,0.20)$. The nonlinear prognosis and covariate-dependent assignment in these designs allow the comparison to extend beyond the location and shape-effect family.

\subsection{Structural-zero designs}\label{sec:zi-design}

The four structural-zero designs use six independent uniform covariates. A unit's outcome vector is zero with probability $1-\pi_a(x)$ and otherwise follows equation~\eqref{eq:dgp} with
\begin{align*}
 m_a^+(x)&=4.80+0.35x_2-0.25x_4+a(\Delta+\beta x_2),\\
 s_a^+(x)&=0.10+0.06x_6-0.03x_2,\qquad
 \gamma_a^+(x)=C_{.05,.85}(0.50+0.06x_4-a\kappa).
\end{align*}
The normal shock standard deviations are $(0.15,0.10)$. The superscript $+$ denotes the nondegenerate branch; Gaussian shocks do not impose a strictly positive support. ZI0 and ZI1 use $(\Delta,\beta,\kappa)=(0,0,0)$, whereas ZI2 and ZI3 use $(0.28,0.08,0.10)$. ZI0 is a placebo and ZI2 changes only the nondegenerate component; both have $\pi_a(x)=C_{.05,.95}\{\sigma(1.2x_4-0.8x_2)\}$. ZI1 changes participation only, adding $a(0.90+0.20x_4)$ to this logistic index. These three designs use IC1's assignment mechanism. ZI3 changes both components and has
\[
 \pi_a(x)=C_{.01,.99}[\sigma\{3(x_4-0.9x_5)+a(0.90+0.60x_4)\}],
 \qquad e(x)=\pi_0(x).
\]
The reference vector is the same as in the IC family.

\subsection{Comparators and implementation}\label{sec:implementation}

DRF estimates the conditional outcome distribution with separate forests for the two arms \citep{DRF-paper}. Causal-DRF uses a shared treatment-aware splitting criterion designed for the conditional kernel treatment effect \citep{naf2026causaldrf}. Both return arm-specific weights on observed training quantile vectors. The comparisons use the authors' implementations, with 2500 trees per arm for DRF and 2500 trees for Causal-DRF, organized in groups of 50 trees. WCF is the particle estimator with functional calibration described in Section~\ref{sec:model}. The structural-zero comparison uses its optional two-part law estimator.

All reported models use the uniform midpoint grid with $K=25$, and WCF uses $M=10$ particles. The boosting budget is 100 iterations, initial step size $0.12$, maximum tree depth 4, minimum total leaf size 10, minimum arm leaf size 5, and smoothing constant $\varepsilon=10^{-3}$. The backtracking search permits twelve halvings and accepts score changes within $10^{-12}$ numerical tolerance. Contrast shrinkage is selected from $\{0,50,500\}$ by observed-arm validation risk, with ties favoring larger penalties. The IC fits use three-fold outcome fitting and selection; the nonlinear and two-part fits use two folds. Propensity estimation uses five-fold logistic regression, with probabilities clipped to $[0.02,0.98]$. This interval excludes part of IC3's true propensity range, so the fitted propensity is not exactly specified in that design even if its untruncated logistic index were known. The mixture probability is fitted with separate logistic classifiers using an $\ell_2$ penalty with inverse strength $C=1$; these classifiers are fitted on the full training arm. Quantile vectors with maximum absolute coordinate at most $10^{-12}$ are classified as degenerate.

The selected outcome penalty is reused for the out-of-fold calibration fits; its selection is not nested within each calibration training fold. Consequently, the out-of-fold predictions exclude each observation from tree fitting but not from penalty selection. The simulations evaluate point-estimation accuracy and do not invoke asymptotic confidence-interval guarantees for this tuning arrangement. The fixed $K$ and $M$ are computational choices; the present comparisons do not establish sensitivity to either resolution.

\subsection{Evaluation}\label{sec:metrics}

We evaluate the mean-quantile contrast $\tau_q(x)$ by the root mean squared error over all test units and all $K$ coordinates. For each scalar functional, TCATE error is the RMSE over the four moderator bins, using the bin average of the true conditional contrast on the test covariates. TATE error is the absolute error in the marginal contrast averaged over those test covariates. For WCF, the recorded marginal predictions average the training-bin AIPW estimates using the test sample's bin frequencies. They therefore evaluate a test-standardized bin estimator, rather than the direct training-sample AIPW mean. Training and test covariates are sampled from the same population. The reference-distance TATE and TCATE use $h_{\mathrm{ref}}$ with the same aggregation rules.

For concise comparisons, functional TATE and TCATE errors are averaged over the four functionals within each replication. Since these functionals have different units, these averages are descriptive summaries whose values depend on the stated outcome scale. Appendix~\ref{app:functional-results} reports each functional separately. Each table gives the mean replication error and its between-replication standard error, after averaging arm-specific or functional-specific entries within the replication. Thus the marginal-effect columns report mean absolute error, not root mean squared error over replications.

Conditional-law error is squared maximum mean discrepancy \citep{gretton2012kernel}, computed with a Gaussian kernel on $(\sqrt{w_k}q_k)_k$. The bandwidth is chosen by a median-distance rule on the quadrature truth nodes, identically across methods within each arm and replication. This bandwidth is used for evaluation only. Law scores average the first 200 test units in each arm and then average the two arms. Oracle expectations use a product Gauss--Hermite rule with twelve nodes per nondegenerate normal shock, separately weighted across mixture components. Accordingly, the law discrepancy is evaluated against the quadrature approximation to the true law.

In the structural-zero designs, mass error is the RMSE between the true zero-component probability and the predicted mass on vectors with maximum absolute coordinate at most $0.05$. The mass-contrast metric is the RMSE of the corresponding treatment contrast over moderator bins. The explicit atom in the two-part model contributes its estimated mixture weight to this probability. All error criteria are oriented so that smaller values indicate greater accuracy.

\section{Results}\label{sec:results}

\subsection{Location and shape-effect designs}

Tables~\ref{tab:revision-income-500} and~\ref{tab:revision-income-1000} report both sample sizes. WCF has the smallest mean error in every displayed column. At $n=1000$, its reference-distance TCATE error is $0.0300$ in IC1, $0.0330$ in IC2, and $0.0463$ in IC3, compared with $0.1333$, $0.1567$, and $0.2476$ for Causal-DRF. Its conditional-law errors also remain lower in these three designs. The IC0 placebo yields smaller false effects for WCF on the mean-quantile, functional, and reference-distance contrasts. These comparisons establish performance differences under the stated assignment and outcome mechanisms; they do not isolate a single architectural cause of the differences.

\begin{table}[t]
\centering
\scriptsize
\setlength{\tabcolsep}{3pt}
\caption{Location and shape-effect designs at $n=500$. Entries are mean(SE) over ten replications. MeanQ is quantile-vector RMSE; Law is squared MMD. TATE columns give mean absolute error and TCATE columns give mean bin RMSE. Functional errors are averaged within each replication.}
\label{tab:revision-income-500}
\begin{tabular}{llrrrrrr}
\toprule
DGP & Method & MeanQ & Law & TATE & TCATE & Ref. TATE & Ref. TCATE \\
\midrule
IC0 & WCF & 0.0503(0.0149) & 0.0205(0.0006) & 0.0143(0.0024) & 0.0384(0.0029) & 0.0245(0.0042) & 0.0496(0.0043) \\
IC0 & Causal-DRF & 0.2323(0.0117) & 0.0651(0.0018) & 0.1371(0.0065) & 0.1374(0.0064) & 0.1869(0.0109) & 0.1870(0.0108) \\
IC0 & DRF & 0.2339(0.0102) & 0.0905(0.0021) & 0.1356(0.0058) & 0.1356(0.0058) & 0.1857(0.0100) & 0.1858(0.0100) \\
IC1 & WCF & 0.1427(0.0116) & 0.0236(0.0010) & 0.0175(0.0022) & 0.0444(0.0025) & 0.0294(0.0059) & 0.0522(0.0051) \\
IC1 & Causal-DRF & 0.2109(0.0114) & 0.0627(0.0023) & 0.1157(0.0061) & 0.1162(0.0061) & 0.1563(0.0102) & 0.1566(0.0101) \\
IC1 & DRF & 0.2194(0.0101) & 0.0821(0.0022) & 0.1188(0.0056) & 0.1189(0.0056) & 0.1611(0.0098) & 0.1612(0.0097) \\
IC2 & WCF & 0.0987(0.0173) & 0.0220(0.0011) & 0.0115(0.0022) & 0.0360(0.0034) & 0.0185(0.0027) & 0.0387(0.0049) \\
IC2 & Causal-DRF & 0.2328(0.0115) & 0.0685(0.0016) & 0.1420(0.0061) & 0.1423(0.0061) & 0.1860(0.0107) & 0.1862(0.0106) \\
IC2 & DRF & 0.2495(0.0127) & 0.0978(0.0023) & 0.1520(0.0067) & 0.1520(0.0067) & 0.2016(0.0121) & 0.2016(0.0121) \\
IC3 & WCF & 0.1415(0.0078) & 0.0251(0.0010) & 0.0290(0.0046) & 0.0550(0.0033) & 0.0285(0.0064) & 0.0668(0.0066) \\
IC3 & Causal-DRF & 0.3444(0.0113) & 0.0773(0.0019) & 0.1950(0.0063) & 0.1954(0.0063) & 0.2662(0.0110) & 0.2666(0.0109) \\
IC3 & DRF & 0.3688(0.0107) & 0.0958(0.0029) & 0.2082(0.0057) & 0.2083(0.0057) & 0.2834(0.0101) & 0.2836(0.0101) \\
\bottomrule
\end{tabular}
\end{table}

\begin{table}[t]
\centering
\scriptsize
\setlength{\tabcolsep}{3pt}
\caption{Location and shape-effect designs at $n=1000$. Entries are mean(SE) over ten replications, using the same metrics as Table~\ref{tab:revision-income-500}.}
\label{tab:revision-income-1000}
\begin{tabular}{llrrrrrr}
\toprule
DGP & Method & MeanQ & Law & TATE & TCATE & Ref. TATE & Ref. TCATE \\
\midrule
IC0 & WCF & 0.0356(0.0054) & 0.0153(0.0005) & 0.0100(0.0012) & 0.0240(0.0025) & 0.0142(0.0025) & 0.0298(0.0028) \\
IC0 & Causal-DRF & 0.1947(0.0081) & 0.0467(0.0010) & 0.1134(0.0048) & 0.1138(0.0047) & 0.1570(0.0081) & 0.1572(0.0080) \\
IC0 & DRF & 0.1963(0.0063) & 0.0641(0.0013) & 0.1104(0.0038) & 0.1105(0.0038) & 0.1543(0.0065) & 0.1544(0.0065) \\
IC1 & WCF & 0.0904(0.0061) & 0.0174(0.0006) & 0.0128(0.0015) & 0.0271(0.0020) & 0.0143(0.0028) & 0.0300(0.0032) \\
IC1 & Causal-DRF & 0.1760(0.0090) & 0.0455(0.0009) & 0.0962(0.0051) & 0.0966(0.0051) & 0.1329(0.0080) & 0.1333(0.0080) \\
IC1 & DRF & 0.1876(0.0058) & 0.0591(0.0012) & 0.1012(0.0032) & 0.1013(0.0032) & 0.1394(0.0056) & 0.1395(0.0056) \\
IC2 & WCF & 0.0700(0.0034) & 0.0167(0.0006) & 0.0099(0.0016) & 0.0255(0.0016) & 0.0143(0.0033) & 0.0330(0.0035) \\
IC2 & Causal-DRF & 0.1975(0.0072) & 0.0488(0.0010) & 0.1176(0.0039) & 0.1179(0.0039) & 0.1565(0.0069) & 0.1567(0.0069) \\
IC2 & DRF & 0.2298(0.0063) & 0.0724(0.0016) & 0.1357(0.0034) & 0.1358(0.0034) & 0.1795(0.0050) & 0.1796(0.0050) \\
IC3 & WCF & 0.1103(0.0043) & 0.0180(0.0006) & 0.0188(0.0027) & 0.0412(0.0030) & 0.0161(0.0042) & 0.0463(0.0064) \\
IC3 & Causal-DRF & 0.3177(0.0076) & 0.0605(0.0013) & 0.1784(0.0043) & 0.1786(0.0042) & 0.2474(0.0067) & 0.2476(0.0067) \\
IC3 & DRF & 0.3481(0.0055) & 0.0773(0.0014) & 0.1972(0.0033) & 0.1973(0.0033) & 0.2701(0.0051) & 0.2702(0.0051) \\
\bottomrule
\end{tabular}
\end{table}


\subsection{Symmetric and nonlinear designs}

Table~\ref{tab:revision-abstract} shows a more varied ranking. WCF has the smallest mean conditional-law error in D0, D2, D5, D7, and D8. In contrast, both forest baselines have smaller reference-distance TCATE error in D2, D5, D6, D7, and D8. Thus improved law approximation under the kernel metric does not ensure improved estimation of every scalar functional. D5 illustrates the distinction: the two potential outcomes have identical mean quantile curves but different conditional laws, so the causal questions depend on the functional chosen.

The multimodal D6 design is particularly challenging for WCF. Its law error is $0.0178$, compared with $0.0064$ for Causal-DRF and $0.0059$ for DRF. Its reference-distance TCATE error is also larger. This result limits claims of general superiority under covariate-dependent assignment and motivates further investigation of the particle representation in multimodal settings.

\begin{table}[t]
\centering
\scriptsize
\setlength{\tabcolsep}{3pt}
\caption{Symmetric and nonlinear designs at $n=1000$. Entries are mean(SE) over ten replications, using the same metrics as Table~\ref{tab:revision-income-500}.}
\label{tab:revision-abstract}
\begin{tabular}{llrrrrrr}
\toprule
DGP & Method & MeanQ & Law & TATE & TCATE & Ref. TATE & Ref. TCATE \\
\midrule
D0 & WCF & 0.0868(0.0021) & 0.0234(0.0015) & 0.0041(0.0005) & 0.0101(0.0005) & 0.0038(0.0006) & 0.0097(0.0008) \\
D0 & Causal-DRF & 0.1391(0.0012) & 0.1033(0.0036) & 0.0230(0.0017) & 0.0512(0.0013) & 0.0087(0.0012) & 0.0886(0.0023) \\
D0 & DRF & 0.1672(0.0022) & 0.1191(0.0041) & 0.0131(0.0010) & 0.0415(0.0012) & 0.0048(0.0016) & 0.0799(0.0023) \\
D2 & WCF & 0.0410(0.0057) & 0.0098(0.0003) & 0.0137(0.0022) & 0.0300(0.0024) & 0.0115(0.0037) & 0.0399(0.0032) \\
D2 & Causal-DRF & 0.0818(0.0052) & 0.0154(0.0003) & 0.0292(0.0043) & 0.0340(0.0039) & 0.0151(0.0053) & 0.0227(0.0048) \\
D2 & DRF & 0.0627(0.0028) & 0.0138(0.0003) & 0.0130(0.0024) & 0.0230(0.0025) & 0.0140(0.0053) & 0.0224(0.0052) \\
D5 & WCF & 0.0946(0.0153) & 0.0220(0.0008) & 0.0146(0.0039) & 0.0332(0.0024) & 0.0202(0.0040) & 0.0418(0.0055) \\
D5 & Causal-DRF & 0.0765(0.0032) & 0.0374(0.0008) & 0.0187(0.0038) & 0.0302(0.0027) & 0.0144(0.0034) & 0.0239(0.0025) \\
D5 & DRF & 0.0917(0.0036) & 0.0431(0.0011) & 0.0132(0.0029) & 0.0295(0.0030) & 0.0151(0.0037) & 0.0235(0.0028) \\
D6 & WCF & 0.2189(0.0178) & 0.0178(0.0008) & 0.0506(0.0127) & 0.0970(0.0147) & 0.0191(0.0063) & 0.0862(0.0094) \\
D6 & Causal-DRF & 0.2310(0.0132) & 0.0064(0.0004) & 0.0504(0.0090) & 0.0845(0.0084) & 0.0261(0.0075) & 0.0585(0.0091) \\
D6 & DRF & 0.1986(0.0099) & 0.0059(0.0003) & 0.0461(0.0077) & 0.0852(0.0052) & 0.0220(0.0060) & 0.0501(0.0068) \\
D7 & WCF & 0.0553(0.0030) & 0.0094(0.0003) & 0.0162(0.0021) & 0.0368(0.0024) & 0.0156(0.0045) & 0.0386(0.0034) \\
D7 & Causal-DRF & 0.0721(0.0034) & 0.0129(0.0003) & 0.0208(0.0039) & 0.0360(0.0032) & 0.0128(0.0042) & 0.0235(0.0033) \\
D7 & DRF & 0.0662(0.0031) & 0.0142(0.0003) & 0.0156(0.0027) & 0.0353(0.0027) & 0.0135(0.0043) & 0.0238(0.0037) \\
D8 & WCF & 0.0635(0.0020) & 0.0127(0.0003) & 0.0234(0.0044) & 0.0672(0.0070) & 0.0377(0.0098) & 0.0793(0.0117) \\
D8 & Causal-DRF & 0.2100(0.0082) & 0.0417(0.0007) & 0.0706(0.0072) & 0.0889(0.0047) & 0.0325(0.0088) & 0.0584(0.0077) \\
D8 & DRF & 0.1623(0.0056) & 0.0416(0.0008) & 0.0257(0.0033) & 0.0509(0.0039) & 0.0446(0.0113) & 0.0724(0.0081) \\
\bottomrule
\end{tabular}
\end{table}


\subsection{Structural-zero outcomes}

Table~\ref{tab:revision-zi} compares the two-part WCF law estimator with the two forest baselines. The two-part estimator has the smallest mean mass-placement and law errors across the four designs. In ZI3, its mass error is $0.0598$, compared with $0.1731$ for Causal-DRF and $0.3475$ for DRF; its law error is $0.0176$, compared with $0.0768$ and $0.2472$. Its mass-contrast errors are also lower in the reported comparison. These results concern estimation of the structural mixture and its participation contrast. The two-part fits do not include the functional AIPW calibration, and no double-robustness claim is attached to these mass or law scores.

\begin{table}[t]
\centering
\scriptsize
\setlength{\tabcolsep}{3pt}
\caption{Structural-zero outcomes at $n=1000$. Entries are mean(SE) over ten replications. Zero mass and mass contrast are RMSEs; Law is squared MMD. Two-part WCF estimates the mixture without AIPW functional calibration.}
\label{tab:revision-zi}
\begin{tabular}{llrrr}
\toprule
DGP & Method & Zero mass & Mass contrast & Law \\
\midrule
ZI0 & Two-part WCF & 0.0537(0.0032) & 0.0427(0.0092) & 0.0133(0.0007) \\
ZI0 & Causal-DRF & 0.0945(0.0026) & 0.0595(0.0054) & 0.0247(0.0013) \\
ZI0 & DRF & 0.1058(0.0029) & 0.0561(0.0079) & 0.0296(0.0013) \\
ZI1 & Two-part WCF & 0.0526(0.0033) & 0.0365(0.0055) & 0.0154(0.0011) \\
ZI1 & Causal-DRF & 0.0901(0.0029) & 0.0426(0.0087) & 0.0249(0.0012) \\
ZI1 & DRF & 0.1022(0.0037) & 0.0479(0.0080) & 0.0291(0.0016) \\
ZI2 & Two-part WCF & 0.0537(0.0032) & 0.0427(0.0092) & 0.0136(0.0007) \\
ZI2 & Causal-DRF & 0.0946(0.0026) & 0.0597(0.0051) & 0.0258(0.0014) \\
ZI2 & DRF & 0.1060(0.0030) & 0.0564(0.0080) & 0.0304(0.0013) \\
ZI3 & Two-part WCF & 0.0598(0.0032) & 0.0445(0.0074) & 0.0176(0.0010) \\
ZI3 & Causal-DRF & 0.1731(0.0039) & 0.2051(0.0082) & 0.0768(0.0028) \\
ZI3 & DRF & 0.3475(0.0079) & 0.3823(0.0130) & 0.2472(0.0099) \\
\bottomrule
\end{tabular}
\end{table}

